\section{Platform Architecture}

\label{sec:architecture}
% ══════════════════════════════════════════════════════════════════════════════

\subsection{System Overview}

ZooPlatform consists of five core subsystems:

\begin{enumerate}
  \item \textbf{Observation Interface}: Mobile and web apps for
    submitting observations with photos, audio, GPS, and metadata.
  \item \textbf{AI Identification}: TaxoNet-powered species suggestions
    with calibrated confidence scores.
  \item \textbf{Gamification Engine}: RL-based personalized challenge
    generation and reward allocation.
  \item \textbf{Quality Model}: Bayesian reliability estimation for
    each observation.
  \item \textbf{Incentive Layer}: ZooRep token distribution for
    verified contributions.
\end{enumerate}

\subsection{Observation Protocol}

ZooPlatform supports two observation modes:

\paragraph{Opportunistic mode.} Any observation, any time. Low barrier
to entry. Used for casual observations and rare species alerts.

\paragraph{Survey mode.} Structured protocol with defined spatial extent,
time duration, and completeness declaration. Produces higher-quality data
suitable for occupancy modeling and abundance estimation.

Both modes capture: species identification, photograph(s), GPS
coordinates ($\pm$5\,m accuracy), timestamp, habitat description, and
optional notes.


% ══════════════════════════════════════════════════════════════════════════════
